\chapter{Metodo delle correnti di maglia}

\section{Cenni teorici}

Il metodo delle correnti di maglia è una tecnica di analisi dei circuiti elettrici lineari in regime stazionario che si basa sull'applicazione della \textbf{Legge di Kirchhoff delle tensioni (LKT)}.  
L'idea fondamentale consiste nell'associare a ciascuna maglia del circuito una corrente fittizia, detta \textbf{corrente di maglia}, e utilizzare tali correnti come incognite del problema.\\

Si definisce \textbf{maglia} un percorso all'interno del circuito con le seguenti caratteristiche:
\begin{itemize}
    \item \textbf{Chiuso}: il percorso deve partire e terminare nello stesso punto.
    \item \textbf{Non intrecciato}: il percorso non deve attraversare uno stesso nodo o uno stesso ramo più di una volta.
\end{itemize}

A ogni maglia viene assegnata una \textbf{corrente di maglia}, il cui verso può essere scelto arbitrariamente; per comodità, si adotta spesso un verso orario per tutte le maglie. È importante sottolineare che le correnti di maglia non coincidono necessariamente con le correnti reali nei rami del circuito: esse rappresentano uno strumento matematico utile per semplificare l'analisi.

Il metodo si fonda sulla \textbf{Legge di Kirchhoff delle tensioni}, secondo la quale la somma algebrica delle tensioni lungo un percorso chiuso è nulla. Tale legge, espressione del principio di conservazione dell'energia, si scrive:

\begin{equation}
\sum_{i} V_i = 0
\label{eq:LKT}
\end{equation}

Per esprimere le tensioni ai capi dei componenti resistivi si utilizza la \textbf{legge di Ohm}, che lega tensione, corrente e resistenza:

\begin{equation}
V = R \cdot I
\label{eq:ohm}
\end{equation}

Nel contesto del metodo delle correnti di maglia, la corrente $I$ che compare nella legge di Ohm può coincidere con una singola corrente di maglia oppure, nel caso di elementi condivisi tra due maglie, con una combinazione algebrica di più correnti.

Un aspetto particolarmente importante riguarda le \textbf{resistenze condivise tra due maglie}. In questo caso, la corrente che attraversa la resistenza non è una sola corrente di maglia, ma la somma o la differenza tra le correnti associate alle due maglie adiacenti. Se, ad esempio, una resistenza è attraversata dalle correnti di maglia $I_1$ e $I_2$, la corrente effettiva nel ramo si può esprimere come:

\begin{equation}
I_R = I_1 \pm I_2
\label{eq:res_condivisa}
\end{equation}

Di conseguenza, la tensione ai capi della resistenza risulta:

\begin{equation}
V_R = R \cdot (I_1 \pm I_2)
\label{eq:tensione_condivisa}
\end{equation}

Il segno tra le due correnti dipende dai versi assegnati alle correnti di maglia e dal verso di percorrenza scelto durante l'applicazione della KVL.

Il metodo delle correnti di maglia è applicabile a circuiti \textbf{lineari e planari} in regime stazionario. Nel caso di circuiti in corrente continua, il regime stazionario implica che le grandezze elettriche non variano nel tempo. In tali condizioni, eventuali condensatori si comportano come circuiti aperti, mentre gli induttori si comportano come cortocircuiti, permettendo di ridurre il circuito a una rete puramente resistiva.

\begin{center}
\fbox{
\parbox{0.9\textwidth}{
\textbf{Osservazione:} le correnti di maglia sono grandezze fittizie introdotte per semplificare la scrittura delle equazioni del circuito.  

Le correnti reali nei rami devono essere determinate a partire dalle correnti di maglia, tenendo conto delle interazioni tra le diverse maglie e dei versi scelti.
}
}
\end{center}



\section{Procedura di risoluzione}

Il metodo delle correnti di maglia si configura come una procedura algoritmica che consente di ridurre un sistema circuitale complesso a un sistema lineare di equazioni. La sua efficacia risiede nella capacità di sfruttare la struttura topologica del circuito per minimizzare il numero di incognite.

Il primo passo consiste in un'analisi topologica del circuito, verificando che si tratti di un circuito planare, ovvero un circuito i cui rami possono essere disegnati su un piano senza intersecarsi se non nei nodi. Successivamente, si identificano le maglie indipendenti, definendo come \textbf{maglia indipendente} una maglia che non può essere ottenuta come combinazione lineare di altre maglie. Il numero di maglie indipendenti \(M\) è dato dalla relazione di Eulero per i grafi planari:
\begin{equation}
M = R - N + 1
\label{eq:topologia}
\end{equation}
dove \(R\) è il numero di rami e \(N\) il numero di nodi. Queste \(M\) maglie costituiscono un insieme sufficiente e necessario per scrivere un sistema di equazioni linearmente indipendenti. Dal punto di vista geometrico, in un circuito planare queste maglie indipendenti corrispondono alle \textbf{maglie elementari}, ossia ai "buchi" che si osservano nel disegno del circuito. Ciascuna maglia elementare rappresenta una regione chiusa del piano che non contiene altre maglie al suo interno; il loro numero è quindi immediatamente visibile dall'aspetto del circuito.

Ad ogni maglia indipendente si attribuisce una variabile, detta \textbf{corrente di maglia}, generalmente indicata con \(I_1, I_2, \dots, I_M\). Il verso di percorrenza di queste correnti è arbitrario; tuttavia, per garantire uniformità e ridurre la probabilità di errori algebrici, si adotta sistematicamente il \textbf{verso orario} per tutte le maglie. Questa scelta, seppur non obbligatoria, semplifica la determinazione dei segni nelle resistenze condivise.

\subsection*{Esempio applicativo}
Per illustrare la procedura, si consideri il circuito di Figura 1, caratterizzato da tre maglie elementari. In questo circuito sono presenti resistori e generatori di tensione indipendenti.


\begin{figure}[h]
\centering
\begin{circuitikz}[scale=1.2, every node/.style={font=\small}]
    % Paths, nodes and wires:
    \draw (0.75, 2) to[american resistor, l_={$R_2$}] (0.75, 4);
    \draw (2.25, 2) to[american resistor, l={$R_3$}] (4, 2);
    \draw (4, 4) to[american resistor, l={$R_4$}] (4, 2);
    \draw (0.75, -0) to[american resistor, l={$R_1$}] (4, -0);
    \draw (0.75, 2) -- (0.75, -0);
    \draw (6, 2) to[american resistor, l_={$R_5$}] (6, 0);
    \draw (4, -0) -- (6, -0);
    \draw (4, 4) -- (6, 4);
    \draw (0.75, 4) to[american voltage source, l={$E_2$}] (4, 4);
    \draw (2.25, 2) to[american voltage source, l={$E_1$}] (0.75, 2);
    \draw (4, -0) to[american voltage source, l={$E_3$}] (4, 2);
    \draw (6, 4) to[american voltage source, l_={$E_4$}] (6, 2);
    % Etichette dei nodi (opzionali, per chiarezza)

    \node[circ] at (4, 4){};	
    \node at (4,4) [above] {A};
	\node[circ] at (4, 0){};
    \node at (4,0) [below] {B};
    \node[circ] at (0.75, 2){};
    \node at (0.75,2) [left] {C};
    \node[circ] at (4, 2){};
    \node at (4,2) [right] {D};
\end{circuitikz}
\caption{Circuito Esempio}
\label{fig:circuito_esempio}
\end{figure}


Cerchiamo ora di definire un procendimento sistematico per ricavare il sistema di equazioni che consente di determinare le correnti di maglia.

\begin{enumerate}
    \item \textbf{Identificare le correnti di maglia:} Assegnare una corrente di maglia a ciascuna maglia indipendente del circuito, scegliendo un verso di percorrenza.
    \item \textbf{Scrivere le equazioni di maglia:} Per ciascuna maglia elementare, applicare la LKT.
    \item \textbf{Esplicitare le tensioni:} Utilizzare la legge di Ohm per esprimere le tensioni ai capi dei resistori in funzione delle correnti di maglia, tenendo conto delle resistenze condivise.
    \item \textbf{BONUS:} Organizzare le equazioni in forma matriciale, identificando la matrice delle resistenze e il vettore dei temrini noti, per applicare metodi risolutivi più efficienti.
\end{enumerate}

\subsection*{1. Identificare le correnti di maglia}
Nel circuito in Figura \ref{fig:circuito_esempio} sono presenti tre maglie elementari. Decidiamo quindi di assegnare 3 correnti di maglia, \(I_1\), \(I_2\) e \(I_3\), tutte con verso orario. Oltre a notare ad occhio che sono presenti 3 maglie elementari, possiamo anche verificare tramite la \ref{eq:topologia}: il circuito ha \(R = 6\) rami e \(N = 4\) nodi, quindi \(M = R-N+1 = 3\) maglie indipendenti, confermando la nostra osservazione.


\begin{figure}[h]
\centering
\begin{circuitikz}[scale=1.2, every node/.style={font=\small}]
    % Paths, nodes and wires:
    \draw (0.75, 2) to[american resistor, l_={$R_2$}] (0.75, 4);
    \draw (2.25, 2) to[american resistor, l={$R_3$}] (4, 2);
    \draw (4, 4) to[american resistor, l={$R_4$}] (4, 2);
    \draw (0.75, -0) to[american resistor, l={$R_1$}] (4, -0);
    \draw (0.75, 2) -- (0.75, -0);
    \draw (6, 2) to[american resistor, l_={$R_5$}] (6, 0);
    \draw (4, -0) -- (6, -0);
    \draw (4, 4) -- (6, 4);
    \draw (0.75, 4) to[american voltage source, l={$E_2$}] (4, 4);
    \draw (2.25, 2) to[american voltage source, l={$E_1$}] (0.75, 2);
    \draw (4, -0) to[american voltage source, l={$E_3$}] (4, 2);
    \draw (6, 4) to[american voltage source, l_={$E_4$}] (6, 2);
    % Etichette dei nodi (opzionali, per chiarezza)

    \node[circ] at (4, 4){};	
    \node at (4,4) [above] {A};
	\node[circ] at (4, 0){};
    \node at (4,0) [below] {B};
    \node[circ] at (0.75, 2){};
    \node at (0.75,2) [left] {C};
    \node[circ] at (4, 2){};
    \node at (4,2) [right] {D};
    % Frecce delle correnti di maglia (aggiunte manualmente)
    \draw[->, thick, red] (1.7,1) arc (120:60:1.2) node[midway, above]{$I_1$};
    \draw[->, thick, blue] (1.7,3) arc (120:60:1.2) node[midway, above]{$I_2$};
    \draw[->, thick, magenta] (5,1) arc (240:110:1) node[midway, right]{$I_3$};
\end{circuitikz}
\caption{Circuito con tre maglie elementari. Le correnti di maglia \(I_1\), \(I_2\) e \(I_3\) sono in senso orario.}
\label{fig:circuito_correnti_maglia}
\end{figure}

\subsection*{2. Scrivere le equazioni di maglia}

Prima di applicare la LKT, è utile ricordare come vengono scelti i segni per i vari componenti con una semplice regola:
\begin{center}
\fbox{
\parbox{0.9\textwidth}{
\textbf{Per i generatori:} il segno della tensione sarà positivo se si attraversa il generatore dal polo negativo al polo positivo, e negativo altrimenti.\\
\textbf{Per i resistori:} il segno della tensione sarà, per convenzione, sempre negativo, poichè si assume che la corrente di maglia attraversi il resistore nel verso della caduta di tensione.
}
}
\end{center}

Scriviamo quindi la LKT per ogni maglia:
\begin{itemize}
    \item \textbf{Maglia 1:} La prima maglia, identificata da \(I_1\), include due resistori e due generatori di tensione. La LKT quindi sarà:
    \begin{equation}
        E_1 - V_{R_3} + E_3 - V_{R_1} = 0
        \label{eq: prima_maglia}
    \end{equation}
    \item \textbf{Maglia 2:} La seconda maglia, identificata da \(I_2\), include tre resistori e due generatori di tensione. La LKT quindi sarà:
    \begin{equation}
        - E_2 - V_{R_4} - V_{R_3} - E_1 - V_{R_2} = 0
        \label{eq: seconda_maglia}
    \end{equation}
    \item \textbf{Maglia 3:} La terza maglia, identificata da \(I_3\), include due resistori e due generatori di tensione. La LKT quindi sarà:
    \begin{equation}
        - E_3 - V_{R_4} - E_4 - V_{R_5} = 0
        \label{eq: terza_maglia}
    \end{equation}
\end{itemize}


\section{Problema del Generatore di Corrente}

\section{Esempi svolti}

\subsection{Esempio 1}
Circuito con soli generatori di tensione e resistenze.

\subsection{Esempio 2}
Circuito con generatore di corrente e più maglie.